\chapter{Introduction}\label{chap:introduction}

\section{Background}

Emergency response is a chain of decisions. A person recognises that something is wrong and asks for help;
a call-taker or system records what is happening and where; someone judges how urgent it is; a dispatcher
chooses which unit to send; a crew travels, treats and, when needed, delivers the patient to a hospital that
must be ready for them. The quality of the whole chain is limited by its weakest link, and every link is
sensitive to time \cite{aringhieri2017ems}. Decades of operations research on emergency medical services have
shown that where vehicles are positioned and how quickly they are assigned strongly influence the time to
reach a patient \cite{gendreau1997tabu,brotcorne2003ambulance}, but those models assume that a good report
and a good assignment decision are already available. In practice they are the hard part.

In Egypt a person in distress must first decide \emph{whom} to call: the police, the ambulance service and
the fire brigade are reached through separate numbers (122, 123 and 180). The report is then a spoken
description, often given in panic, that a human must interpret, locate and prioritise. Once units are on their
way, the person who asked for help usually knows nothing more, and the hospital that will receive a patient
typically learns about them only on arrival. Each of these gaps is a place where minutes are lost.

At the same time, the software that could close the gaps has become far more accessible. Smartphones provide
precise positioning and cameras; managed cloud services provide identity, storage and messaging; open
geographic data and routing engines are freely available \cite{haklay2008osm,luxen2011osrm}; and large language
models can classify free text at a level that was out of reach a few years ago \cite{brown2020gpt3}. A small team
can now assemble, with modern software engineering practice \cite{pressman2020se,sommerville2016software}, a
platform that would once have required a large organisation.

\section{Problem Statement}\label{sec:problem}

An emergency-response platform for a country like Egypt has to solve several distinct problems at once:

\begin{enumerate}
  \item \textbf{Reporting is fragmented and fragile.} A citizen should be able to ask for help in one action,
        without choosing a service, even with a poor connection, and without the system creating duplicate
        emergencies when a request is retried.
  \item \textbf{Reports arrive unstructured and unprioritised.} Dispatchers need a live, ordered queue in which
        the most urgent reports are visible first and duplicate reports of the same event are recognised.
  \item \textbf{Coordination is opaque.} Dispatchers need to see which units are free and where they are;
        crews need a clear mission workflow; and the reporter needs to know that help is coming.
  \item \textbf{Hospitals are not prepared.} An ambulance crew should be able to choose a receiving hospital
        and send the patient's vital signs ahead, so the hospital can prepare before the patient arrives.
  \item \textbf{Medical information is sensitive but life-saving.} A patient's allergies and conditions help the
        people treating them, but must never be exposed to anyone else.
  \item \textbf{The service can be abused.} False reports waste scarce resources, yet an emergency line must
        never refuse a genuine caller.
\end{enumerate}

Existing solutions tend to address one of these problems in isolation. Citizen-facing applications improve
reporting but stop at the point where a human takes over; dispatch software organises resources but expects
the report to arrive by telephone; and academic prototypes frequently reduce the whole domain to a form that
an administrator approves. No single system connects the citizen, the dispatcher, the field crew and the
hospital in one auditable workflow while treating privacy and abuse as first-class design concerns
(\cref{chap:literature}).

\section{Proposed Solution}

\textbf{\ProjectName} (\ProjectFullName) is a complete emergency-response platform that connects all four
parties in a single workflow. A citizen files a report with one tap. The platform captures the location,
classifies the report's priority and flags possible duplicates with a language model, and places the report in a
live queue. A dispatcher reviews the queue and offers the work to the nearest suitable response unit. The unit's
crew accepts, travels, arrives and completes the mission from a mobile app; an ambulance crew additionally selects
a receiving hospital and streams the patient's vital signs to it during transport. Throughout, the reporter
follows every step live and can talk to the people responding.

The platform serves ten roles --- citizen, first responder, ambulance crew, police officer, firefighter,
dispatcher, hospital staff, facility manager, administrator and super administrator --- through two clients and one
backend (\cref{fig:intro-context}):

\begin{itemize}
  \item a \textbf{Flutter mobile app} for citizens and field responders;
  \item a \textbf{Next.js web console} for every role, and the main workspace of dispatchers, hospital staff,
        facility managers and administrators;
  \item a \textbf{Spring Boot API} that enforces every business rule and persists to PostgreSQL, supported by a
        separately deployable \textbf{AI model service} and by managed services for identity, storage, messaging,
        identity verification and maps.
\end{itemize}

\diagram{arch_context.png}{\ProjectName{} at a glance: two clients, one backend and the external services it relies on.}{fig:intro-context}

Artificial intelligence plays a deliberately limited role. It \emph{recommends}: it suggests a priority and a possible
duplicate. A human always decides, and a failure of the AI never affects the handling of a report.

\section{Project Objectives}

The primary objective is to design and implement a platform that shortens the path from an emergency to a
properly equipped response, while remaining safe, private and trustworthy. The specific objectives are to:

\begin{enumerate}
  \item shorten the time from an emergency to a unit being offered the work, by giving dispatchers a prioritised,
        de-duplicated, live queue;
  \item make reporting safe and fast under stress: one-tap SOS, automatic location, and correct behaviour through
        poor connectivity;
  \item give every actor exactly the visibility their role requires, including a consent-based, audited path to a
        patient's medical profile;
  \item keep the service trustworthy through identity verification, a graduated caller-trust policy that never
        blocks a genuine emergency, and an audit trail for sensitive actions;
  \item coordinate the hand-off to hospitals, with live vital signs during transport; and
  \item operate in Arabic and English, in both text directions, on phones and desktops.
\end{enumerate}

\section{Project Scope}

\subsection*{In scope}
\begin{itemize}
  \item Account lifecycle: e-mail and Google registration, e-mail verification by one-time code, phone verification,
        profile completion and administrator-provisioned staff accounts.
  \item Identity verification, versioned legal consent, and a citizen medical profile with consent controls.
  \item Emergency reporting for three categories (medical, police, fire), evidence attachments, live tracking,
        chat, history and first-aid guidance.
  \item Dispatch operations: live queue, AI priority and duplicate suggestion, unit assignment, cancellation,
        false-report handling and caller history.
  \item Field operations: shifts with a unit lead and crew, mission lifecycle, live unit location, hospital
        selection and patient-vitals streaming.
  \item Administration: facilities, response units, users and roles, responder applications, guidance content, legal
        documents; and facility management for one facility's units and staff.
  \item Notifications (in-app inbox with push delivery) and platform services (status, update check, support form).
\end{itemize}

\subsection*{Out of scope}
\begin{itemize}
  \item Direct integration with national emergency telephone numbers or legacy computer-aided dispatch systems.
        \ProjectName{} is designed to operate alongside existing services, not to replace or connect to them.
  \item Payment, billing or insurance processing.
  \item Medical diagnosis or treatment recommendations: the platform transports information; it does not advise on care.
  \item An iOS build of the mobile app. The code base is cross-platform, but only Android is built and distributed.
\end{itemize}

\section{Development Methodology}

\ProjectName{} was developed iteratively and incrementally, following the principles of agile development
\cite{beck2001agile}: each feature is taken in a small vertical slice through database migration, backend rule,
API, both clients and tests before the next one begins. The database schema is versioned with Flyway and validated
against the entities when the backend starts, so a code change that disagrees with the schema fails at once rather
than in production. Continuous integration builds and tests the mobile app and the AI service on every change and
produces a signed Android release when a version tag is pushed. The rules that matter most --- caller trust,
identity-verification policy, webhook verification, one-time-code handling --- are kept in small, pure, unit-tested
classes (\cref{chap:implementation}).

\section{Contributions}

The main contributions of this work are:

\begin{itemize}
  \item a \textbf{complete, working platform} that connects citizens, dispatchers, field crews and hospitals, built
        from a Flutter app, a Next.js console, a Spring Boot API and an AI model service;
  \item a set of \textbf{reliability patterns} suited to emergency software: idempotent submission, an offline outbox,
        after-commit side effects and graceful degradation when an external service fails;
  \item a \textbf{privacy-by-design medical-profile model}: standing consent plus per-report choice, access restricted to
        treating roles and open incidents, every read audited and failing closed;
  \item a \textbf{graduated caller-trust policy} that slows repeated abuse without ever refusing a genuine emergency;
  \item a \textbf{versioned AI model-serving service} with a strict input contract, ready to host clinical models; and
  \item \textbf{complete engineering documentation}: requirements, 64 use cases, 41 diagrams, a 26-table data
        dictionary, an API reference and a traceability matrix.
\end{itemize}

\section{Organisation of the Book}

\Cref{chap:requirements} presents the stakeholders, the functional and non-functional requirements, the user
stories, the use cases and the business rules. \Cref{chap:literature} reviews the literature on emergency
dispatch, decision support, software architecture, real-time communication and privacy, and identifies the gap
that \ProjectName{} fills. \Cref{chap:design} describes the design: architecture, data model, classes, states,
behaviour, security and deployment. \Cref{chap:implementation} explains how the design was implemented, with
excerpts of the real code. \Cref{chap:testing} describes the verification strategy and the scenarios used to check
the requirements. \Cref{chap:conclusion} concludes with an assessment of the objectives, the limitations and the
future work. The appendices contain the database reference, the REST API reference, further screens, the complete
use-case specifications, the design decisions and the configuration reference.
